%!TEX root = ../rigorous_dyson_theorem.tex

\section{Proofs of Main Theorem Components (Section 3)}

This appendix contains the proofs of lemmas and propositions from Section 3, which together establish the main theorem.

\subsection{Proof of Lemma \ref{lem:stein}}
% Proof of Lemma \ref{lem:stein}
\begin{proof}
This is the standard Gaussian integration by parts (Stein's lemma). For $X_{ij} \sim \mathcal{N}(0, \sigma^2/N)$ independent of other variables, and any smooth function $\phi$:
\begin{align}
\E[X_{ij} \phi(X)]
&= \int X_{ij} \phi(X) \cdot \frac{\sqrt{N}}{\sigma\sqrt{2\pi}} e^{-N X_{ij}^2/(2\sigma^2)} dX_{ij} \\
&= \int \phi(X) \cdot \Bigl(-\frac{\sigma^2}{N} \frac{\partial}{\partial X_{ij}}\Bigr) \frac{\sqrt{N}}{\sigma\sqrt{2\pi}} e^{-NX_{ij}^2/(2\sigma^2)} dX_{ij} \\
&= \frac{\sigma^2}{N} \int \frac{\partial \phi}{\partial X_{ij}}(X) \cdot \frac{\sqrt{N}}{\sigma\sqrt{2\pi}} e^{-NX_{ij}^2/(2\sigma^2)} dX_{ij} \\
&= \frac{\sigma^2}{N} \E\Bigl[\frac{\partial \phi}{\partial X_{ij}}\Bigr].
\end{align}
Applying componentwise to $F$ gives the result.
\end{proof}


\subsection{Proof of Proposition \ref{prop:self-energy}}
% Proof of Proposition \ref{prop:self-energy}
\begin{proof}
We compute each block of $\mcX \mcR$ separately.

\textbf{Step 1: Block structure.} We have
\[
\mcX = \begin{pmatrix} 0 & X \\ X^\top & 0 \end{pmatrix}, \quad
\mcR = \begin{pmatrix} R_{11} & R_{12} \\ R_{21} & R_{22} \end{pmatrix},
\]
so
\[
\mcX \mcR = \begin{pmatrix} X R_{21} & X R_{22} \\ X^\top R_{11} & X^\top R_{12} \end{pmatrix}.
\]

\textbf{Step 2: Isotropy-based evaluation.} We need to evaluate $\E[\mcX \mcR]$ in terms of $\E[\mcR]$. Since $\E[X] = 0$, the naive expansion $\E[\mcX]\E[\mcR] = 0$ vanishes. The leading non-zero contribution comes from Gaussian integration by parts combined with the isotropy condition.

From Step 1, we have
\[
\mcX \mcR = \begin{pmatrix} X R_{21} & X R_{22} \\ X^\top R_{11} & X^\top R_{12} \end{pmatrix}.
\]

The key insight is that the self-energy $\E[\mcX \mcR]$ involves quadratic forms in $X$ which can be evaluated directly using the isotropy condition.

\textbf{Step 3: Apply isotropy to key contractions.} The dominant contribution comes from terms of the form $\E[X M X^\top]$ for deterministic or weakly random matrices $M$. By isotropy (Definition \ref{def:isotropy}),
\begin{equation}\label{eq:isotropy-contraction-1}
\boxed{\E[X M_{21} X^\top] = \frac{\sigma^2}{N}\Tr(M_{21}) I}
\end{equation}
for any deterministic matrix $M_{21}$.

Similarly, for the $(2,1)$ block of $\E[\mcX \mcR]$, which is $\E[X^\top R_{11}]$, we have
\begin{equation}\label{eq:isotropy-contraction-2}
\boxed{\E[X^\top M_{12} X] = \frac{\sigma^2}{N}\Tr(M_{12}) I.}
\end{equation}

\textbf{Step 4: Leading-order self-energy.} Applying \eqref{eq:isotropy-contraction-1} and \eqref{eq:isotropy-contraction-2} with $M = \E[\mcR]$, we obtain
\begin{align}
\E[(\mcX \mcR)_{12}] &= \E[X R_{22}] \\
&\approx \sigma^2 \cdot \frac{1}{N}\Tr(\E[R_{21}]) \cdot I + \text{(correction terms)} \\
&= \sigma^2 m_{21}^{(0)} I + O(1/N),
\end{align}
where $m_{21}^{(0)} = \frac{1}{N}\Tr(\E[R_{21}])$ is the zeroth-order approximation.

More carefully, the Wick contraction via Gaussian integration by parts (using Lemma \ref{lem:stein} and the isotropy condition) gives:
\begin{align}
\E[\mcX \mcR]
&\approx \begin{pmatrix} 0 & \sigma^2 \frac{1}{N}\Tr(\E[R_{12}]) I \\ \sigma^2 \frac{1}{N}\Tr(\E[R_{21}]) I & 0 \end{pmatrix} \cdot \E[\mcR] + \text{Error}.
\end{align}

Define
\begin{equation}
\Sigma[\mcM] := \begin{pmatrix} 0 & \sigma^2 m_{12} I \\ \sigma^2 m_{21} I & 0 \end{pmatrix},
\quad
m_{ij} = \frac{1}{N}\Tr(M_{ij}),
\end{equation}
where $\mcM = \E[\mcR]$.

Then we have established
\begin{equation}\label{eq:wick-closure}
\E[\mcX \mcR] = \Sigma[\E[\mcR]] \cdot \E[\mcR] + \text{Error}.
\end{equation}

\textbf{Step 5: Error bound.} The error terms come from:
\begin{enumerate}
\item \textbf{Replacement error:} Replacing $R_{ij}$ with $\E[R_{ij}]$ inside the isotropy formula. By resolvent bounds (Lemma \ref{lem:resolvent-bounds}), $\norm{R_{ij}} \leq C/\eta$, so
\[
\norm{\E[X (R_{21} - \E[R_{21}]) X^\top]}_F \leq \frac{\sigma^2}{N} \norm{R_{21} - \E[R_{21}]}_F \cdot N \leq C \norm{R_{21} - \E[R_{21}]}_F.
\]
This couples to the error we're bounding (self-consistent equation below).

\item \textbf{Trace approximation error:} Replacing the full matrix $M_{ij}$ with the scalar $\frac{1}{N}\Tr(M_{ij}) \cdot I$. We have
\begin{align}
\norm{M_{ij} - \frac{1}{N}\Tr(M_{ij}) I}_F^2
&= \norm{M_{ij}}_F^2 - \frac{1}{N}|\Tr(M_{ij})|^2 \\
&\leq \norm{M_{ij}}_F^2 \leq \frac{C}{\eta^2} N.
\end{align}
Thus $\norm{M_{ij} - \frac{1}{N}\Tr(M_{ij}) I}_F \leq C\sqrt{N}/\eta$, giving a contribution of order $O(1/\sqrt{N})$ to the error.

\item \textbf{Higher-order Wick contractions:} Terms with four or more $X$ factors. By Wick's theorem, these contribute at most
\[
\E[\norm{X}^4 \norm{R}^2] \leq \frac{C}{N^2} \cdot \frac{1}{\eta^2} = O(1/N^2),
\]
using the fourth moment bound (Assumption A4).
\end{enumerate}

Combining these estimates, we obtain
\begin{equation}
\norm{\text{Error}}_F \leq \frac{C(\eta, \sigma, \norm{A})}{N}.
\end{equation}

This completes the proof of equation \eqref{eq:self-energy-formula}.
\end{proof}


\subsection{Proof of Proposition \ref{prop:fixed-point}}
% Proof of Proposition \ref{prop:fixed-point}
\begin{proof}
We prove existence and uniqueness using the Banach contraction mapping theorem on $\C$.

\textbf{Step 1: Define the space.}
Let
\[
\mcX = \bigl\{x \in \C : |x| \leq R\bigr\}
\]
for $R > 0$ to be determined. Equipped with the standard absolute value metric, $\mcX$ is a closed subset of the complete metric space $\C$.

\textbf{Step 2: Explicit formula for $K(x)$.}
For $x \in \mcX$, the perturbed linearization is
\[
\mcL_0(z) - \Sigma(x)
= \begin{pmatrix} zI & -(A + \sigma^2 x I) \\ -(A^\top + \sigma^2 x I) & I \end{pmatrix}.
\]

By Lemma \ref{lem:block-inversion}, the $(1,2)$ block of the inverse is:
\begin{align}
\bigl((\mcL_0 - \Sigma(x))^{-1}\bigr)_{12}
&= -(zI)^{-1}(A + \sigma^2 x I) \bigl(I - (A^\top + \sigma^2 x I)(zI)^{-1}(A + \sigma^2 x I)\bigr)^{-1} \notag\\
&= \frac{1}{z}(A + \sigma^2 x I) \bigl(I - \tfrac{1}{z}(A^\top + \sigma^2 x I)(A + \sigma^2 x I)\bigr)^{-1}. \label{eq:K-explicit}
\end{align}

Taking the normalized block trace:
\begin{equation}\label{eq:K-formula}
K(x) = \frac{1}{Nz}\Tr\left((A + \sigma^2 x I)\left(I - \tfrac{1}{z}(A^\top + \sigma^2 x I)(A + \sigma^2 x I)\right)^{-1}\right).
\end{equation}

\textbf{Step 3: Bounded operator.}
We show $K: \mcX \to \mcX$ for appropriate choice of $R$.

Define $H(x) = (A^\top + \sigma^2 x I)(A + \sigma^2 x I)$. This matrix is self-adjoint and satisfies
\[
\norm{H(x)} \leq (\norm{A} + \sigma^2 |x|)^2 \leq (\norm{A} + \sigma^2 R)^2.
\]

For $|z| \geq \eta$ and $\eta > (\norm{A} + \sigma^2 R)^2/\eta$, the resolvent is bounded:
\[
\norm{\bigl(I - \tfrac{1}{z}H(x)\bigr)^{-1}} \leq \frac{|z|}{|z| - \norm{H(x)}/|z|} \leq \frac{\eta}{\eta - (\norm{A} + \sigma^2 R)^2/\eta}.
\]

Thus:
\begin{align}
|K(x)|
&\leq \frac{1}{|z|} \cdot \frac{1}{N}\norm{A + \sigma^2 x I} \cdot \norm{\bigl(I - \tfrac{1}{z}H(x)\bigr)^{-1}} \cdot \sqrt{N} \notag\\
&\leq \frac{1}{\eta} \cdot (\norm{A} + \sigma^2 R) \cdot \frac{\eta}{\eta - (\norm{A} + \sigma^2 R)^2/\eta}. \label{eq:K-bound}
\end{align}

Choose $R = \frac{2(\norm{A} + 1)}{\eta}$. Then for $\eta^2 > 2(\norm{A} + \sigma^2 R)^2$, we have
\[
|K(x)| \leq \frac{(\norm{A} + \sigma^2 R)}{\eta - (\norm{A} + \sigma^2 R)^2/\eta} \leq R/2 < R.
\]

Thus $K$ maps $\mcX$ to $\mcX$ when $\eta > \sqrt{2}(\norm{A} + 1)\sigma$.

\textbf{Step 4: Lipschitz constant via resolvent derivative.}
For $x^1, x^2 \in \mcX$, define $R_1 = (\mcL_0 - \Sigma(x^1))^{-1}$ and $R_2 = (\mcL_0 - \Sigma(x^2))^{-1}$.

By the resolvent identity (Lemma \ref{lem:resolvent-identity}):
\[
R_1 - R_2 = R_1 \bigl(\Sigma(x^2) - \Sigma(x^1)\bigr) R_2.
\]

The self-energy difference satisfies:
\[
\Sigma(x^2) - \Sigma(x^1) = \sigma^2(x^2 - x^1) \begin{pmatrix} 0 & I \\ I & 0 \end{pmatrix},
\quad
\norm{\Sigma(x^2) - \Sigma(x^1)} \leq 2\sigma^2 |x^2 - x^1|.
\]

Taking the $(1,2)$ block and extracting the block trace:
\begin{align}
|K(x^1) - K(x^2)|
&= \bigl|\bTr\bigl((R_1)_{12} - (R_2)_{12}\bigr)\bigr| \notag\\
&\leq \frac{1}{N}\norm{(R_1)_{12} - (R_2)_{12}}_F \cdot \sqrt{N} \notag\\
&\leq \norm{R_1} \cdot \norm{\Sigma(x^2) - \Sigma(x^1)} \cdot \norm{R_2} \notag\\
&\leq \frac{C^2}{\eta^2} \cdot 2\sigma^2 |x^2 - x^1|, \label{eq:lipschitz-K}
\end{align}
where $C = C(\norm{A})$ is the resolvent bound constant.

For $\eta > \sqrt{2}C\sigma$, the Lipschitz constant $L = 2C^2\sigma^2/\eta^2 < 1$.

\textbf{Step 5: Apply Banach fixed-point theorem.}
Under the condition $\eta > \max\{\sqrt{2}(\norm{A} + 1)\sigma, \sqrt{2}C\sigma\} = 2(1 + \norm{A})\sigma$, the operator $K: \mcX \to \mcX$ is a contraction on the closed subset $\mcX$ of the complete metric space $(\C, |\cdot|)$.

By the Banach fixed-point theorem, there exists a unique fixed point $x^* \in \mcX$ satisfying $x^* = K(x^*)$.

\textbf{Step 6: Herglotz property.}
We verify that $\Im(m_{11}) > 0$ where $m_{11}$ is given by \eqref{eq:m11-from-x}.

Define $H = (A + \sigma^2 x^* I)(A^\top + \sigma^2 x^* I)$. This matrix is self-adjoint and positive semidefinite.

By the spectral theorem, $H = U \Lambda U^\top$ with $U$ orthogonal and $\Lambda = \mathrm{diag}(\lambda_1, \ldots, \lambda_N)$ with $\lambda_j \geq 0$. Thus:
\[
(zI - H)^{-1} = U \mathrm{diag}\Bigl(\frac{1}{z - \lambda_j}\Bigr) U^\top.
\]

For $z = x + iy$ with $y = \Im(z) > 0$ and real $\lambda_j \geq 0$:
\[
\Im\Bigl(\frac{1}{z - \lambda_j}\Bigr) = \frac{y}{|z - \lambda_j|^2} > 0.
\]

By the spectral decomposition:
\[
\Im((zI - H)^{-1}) = U \cdot \mathrm{diag}\Bigl(\frac{y}{|z - \lambda_j|^2}\Bigr) \cdot U^\top \succeq 0.
\]

Since this is positive semidefinite, its normalized trace is positive:
\[
\Im(m_{11}) = \Im\Bigl(\frac{1}{N}\Tr((zI - H)^{-1})\Bigr) = \frac{1}{N}\sum_{j=1}^N \frac{y}{|z - \lambda_j|^2} > 0.
\]

\textbf{Step 7: Dyson equation.}
Setting $m_{12} = m_{21} = x^*$ and computing $m_{11}, m_{22}$ via \eqref{eq:m11-from-x}--\eqref{eq:m22-from-x}, the $2 \times 2$ matrix $m = (m_{ij})$ satisfies $m = T(m)$ where $T$ is the operator from the original matrix formulation.

The corresponding $2N \times 2N$ block matrix $\mcM(z) = \mcM[m]$ with blocks $M_{ij} = m_{ij} I_N$ satisfies the Dyson equation \eqref{eq:dyson-fp}. Continuity of $x(z)$ in $z$ follows from the contraction mapping theorem.
\end{proof}


\subsection{Proof of Lemma \ref{lem:concentration}}
% Proof of Lemma \ref{lem:concentration}
\begin{proof}
We use the Lipschitz concentration inequality for Gaussian random variables.

\textbf{Step 1: Lipschitz constant.} The function $F(X) = \bTr(\mcR(X))$ depends on $X$ through the resolvent. We compute the Lipschitz constant with respect to the Frobenius norm on $X$.

Let $X, X'$ differ in a single entry $(i, j)$ by amount $\Delta$. Then
\begin{align}
\norm{\mcR(X) - \mcR(X')}_F
&\leq \norm{\mcR(X)(\mcL(X') - \mcL(X))\mcR(X')}_F \\
&= \norm{\mcR(X) \cdot \Delta \cdot e_i e_j^\top \cdot \mcR(X')}_F \\
&\leq |\Delta| \cdot \norm{\mcR(X)} \cdot \norm{\mcR(X')} \\
&\leq \frac{|\Delta|}{\eta^2},
\end{align}
using Corollary \ref{cor:block-resolvent-bounds}.

Thus the Lipschitz constant is $L = \frac{1}{\eta^2}$ per entry.

\textbf{Step 2: Total Lipschitz bound.} Summing over all $N^2$ entries with variance $\sigma^2/N$ each:
\[
\text{Total Lipschitz constant} \leq \frac{N^2 \cdot \sigma^2/N}{\eta^4} = \frac{N\sigma^2}{\eta^4}.
\]

\textbf{Step 3: Gaussian concentration.} By the Gaussian concentration inequality (see Appendix B),
\[
\mathbb{P}(|F(X) - \E[F(X)]| > t) \leq 2\exp\Bigl(-\frac{t^2}{2L^2}\Bigr),
\]
where $L^2 = \frac{C \sigma^2}{\eta^4 N}$ is the variance-adjusted Lipschitz constant.

Substituting gives the result.
\end{proof}


\subsection{Proof of Theorem \ref{thm:main}}

The main theorem has four parts, proven as follows:

\textbf{Part A (Deterministic Equivalent):} The existence and uniqueness of the deterministic equivalent $\mcM(z)$ satisfying the Dyson equation follows directly from Proposition \ref{prop:fixed-point}. The scalar fixed-point formulation exploits the symmetry $m_{12} = m_{21}$ (Lemma \ref{lem:symmetry}) to reduce the problem to finding a unique $x \in \C$ satisfying $x = K(x)$, which is proven via Banach contraction mapping theorem.

\textbf{Part B (Almost Sure Convergence):} This is proven in Section 3.5 using the self-consistent error equation. The proof combines Proposition \ref{prop:self-energy} (which establishes the self-energy structure) with Lemma \ref{lem:concentration} (which provides concentration bounds). The convergence follows from showing that the error $E_N = \bTr(\E[\mcR_N]) - \bTr(\mcM)$ satisfies a contractive equation with $O(1/N)$ residual.

\textbf{Part C (Convergence Rate):} The optimal $O(1/N^2)$ convergence rate is established in Section 4 (see Theorem \ref{thm:convergence-rate} and Appendix C).

\textbf{Part D (Spectral Convergence):} This follows from Part B via the Stieltjes inversion formula (Lemma \ref{lem:stieltjes}). The almost sure convergence of $\mcG_N(z)$ to $\bTr(\mcM(z))$ implies weak convergence of the empirical spectral measure via standard arguments from random matrix theory.
